\documentclass{article}
\usepackage{iclr2027_conference,times}
\usepackage[utf8]{inputenc}
\usepackage[T1]{fontenc}
\usepackage{amsmath,amssymb,amsthm}
\usepackage{booktabs,multirow,array}
\usepackage{graphicx,placeins}
\usepackage{xcolor,colortbl}
\usepackage{url,microtype,hyperref}
\hypersetup{hidelinks}

% The official ICLR style is unmodified. Color is confined to figures and tables.
\definecolor{Navy}{HTML}{17324D}
\definecolor{JEPA}{HTML}{CD663E}
\definecolor{Real}{HTML}{087E83}
\definecolor{Flow}{HTML}{7860A5}
\newcommand{\inputtable}[1]{\input{tables/#1}}
\newcommand{\paperfig}[2]{\IfFileExists{../shared/figures/#1}{\includegraphics[width=#2\linewidth,clip]{../shared/figures/#1}}{\fbox{\parbox[c][1.4in][c]{0.95\linewidth}{\centering Figure asset: #1}}}}
\newtheorem{proposition}{Proposition}
\newtheorem{lemma}{Lemma}
\newtheorem{corollary}{Corollary}
\theoremstyle{remark}
\newtheorem{remark}{Remark}
\newcommand{\E}{\mathbb{E}}
\newcommand{\Var}{\mathrm{Var}}
\newcommand{\Cov}{\mathrm{Cov}}
\newcommand{\tr}{\mathrm{tr}}
\newcommand{\R}{\mathbb{R}}
\newcommand{\Nrm}{\mathcal{N}}
\newcommand{\Lcal}{\mathcal{L}}

\title{Learning to Ignore What Matters:\\Representation, Prediction, and Control\\in JEPA World Models}
\author{Anonymous authors\\Paper under double-blind review}

\begin{document}
\maketitle
\lhead{Under review as a conference paper at ICLR 2027}
\begin{abstract}
A world model should discard visual detail without discarding the distinctions an agent needs to act.
Joint-embedding prediction makes this trade-off implicit: it learns both the representation and the
prediction target. We examine three separate questions: what information remains encoded, what
survives prediction, and what supports control. We begin with Craftax, where one 512-dimensional
joint-embedding model reaches 121 of 128 nearby goals in a six-action open-loop test, yet its tested
reward-planning pipeline underperforms random actions. A controlled moving-wall task then separates
representation from imagination: decoded real encodings retain the wall, and policies using them
exceed $50\%$ survival in two of nine runs; decoded predicted rollouts lose the wall, and none of 27 policies
trained in the tested imagination pipeline reaches that threshold. These pipeline outcomes include fitted heads
and policy optimization and do not isolate dynamics as the cause. Finally, nuisance-texture worlds
connect the observations to the objective. We formalize a known predictable-nuisance failure under
Gaussian non-collapse, and construct an exactly Gaussian, fully decodable representation in which
a control-relevant bit has arbitrarily small prediction-loss scale. The resulting lesson is a concrete
evaluation requirement: retention, prediction, and control must be checked separately.
\end{abstract}

\section{Introduction}
\label{sec:intro}

A useful world model need not reproduce every pixel. It must preserve the distinctions that change
what an agent should do. For a robot approaching a doorway, the texture of the wall may be irrelevant;
the wall's position and the available opening are not. Learning this separation without specifying a
particular downstream reward is a central attraction of joint-embedding predictive architectures
(JEPAs): predict an informative representation of the future rather than reconstruct the future
image \citep{lecun2022path,assran2023ijepa,bardes2024vjepa}.

Recent systems make this approach increasingly practical. Pretrained video representations support
action-conditioned planning, and end-to-end objectives jointly learn representations and dynamics
\citep{assran2025vjepa2,sobal2025pldm,maes2026lewm}. These successes raise an evaluation question:
\emph{when a learned world model supports one behavior, what has that established about its usefulness
for another?} Nearby goal reaching, recognizable decoded rollouts, and reward-driven control measure
different capabilities. A representation can encode a feature that its predictor does not preserve;
a predictor can achieve low error on features that a controller cannot use.

We study these distinctions in a squared-loss JEPA with Gaussian marginal regularization, following
the end-to-end formulation of LeWorldModel \citep{maes2026lewm,balestriero2025lejepa}. The study moves
from a complex environment to progressively simpler ones (Figure~\ref{fig:staircase}). Craftax first
exposes a split between short-horizon goal reaching and reward control. A moving-wall task reduces
the decision to one observable obstacle and an initially hidden gap. Cube worlds then remove the
obstacle and vary only visual nuisance and dynamics. The simplification makes the question sharper:
is relevant information absent from the encoding, lost during prediction, or ineffective in the
complete control pipeline?

The gate provides the clearest empirical separation. A decoder recovers the wall from frozen real
encodings, and a convolutional policy trained on those decoded observations discovers the survival
strategy in two of nine runs. The same decoder applied to predicted latents loses the wall within
two steps for all three fitted encoders. None of the 27 policies trained in the tested JEPA
imagination pipeline discovers the strategy. The decoder adds supervision, and imagined training
also changes reward estimation and policy optimization; these distinctions are essential to what
the experiment establishes.

Our contributions are threefold. First, we organize the evidence around separate tests of
\textbf{retention, prediction, and control}, making clear which transfers between these tests are
supported. Second, we provide a controlled, replicated gate study showing useful information in a
frozen representation alongside poor performance of its imagination pipeline, with predictor
variants and model seeds reported separately. Third, we characterize two limits of the training
objective: Gaussian non-collapse can coexist with nuisance-only encoding, and exact decodability
can coexist with arbitrarily small prediction-loss scale. The predictable-nuisance mechanism is
known \citep{sobal2022slow}; our theory specializes it to Gaussian regularization and separates it
from the distinct question of how much loss is assigned to retained information.

\begin{figure}[t]
\centering
\paperfig{evidence_staircase.pdf}{1.0}
\caption{\textbf{Simplify the environment to sharpen the question.} Craftax tests whether local
visual-goal planning transfers to reward-driven behavior. The gate exposes a single control-relevant
obstacle and separates encoded observations from imagined futures. Cube worlds isolate predictable
nuisance. The environments supply complementary diagnostics, not a shared performance leaderboard.}
\label{fig:staircase}
\end{figure}

\section{What is learned, and what is evaluated?}
\label{sec:worlds}

\paragraph{Joint representation and prediction.}
An encoder maps an image to $z_t=f_\theta(x_t)\in\R^D$. An action-conditioned predictor uses
$H$ latent frames and their actions to predict the next latent:
\begin{equation}
\label{eq:lewm}
\Lcal(\theta)=\frac{1}{D}\widehat{\E}_t\left\|g_\theta(z_{t-H+1:t},a_{t-H+1:t})-z_{t+1}\right\|^2
+\lambda\,\mathrm{SIGReg}(\{z_t\}).
\end{equation}
Gradients reach both the predictions and their targets; there is no stop-gradient or teacher.
SIGReg compares the characteristic functions of random latent projections with those of a standard
Gaussian \citep{balestriero2025lejepa}. It constrains the marginal distribution of features, not which
physical variables they represent. The precise empirical and population expressions appear in
Appendix~\ref{app:proofs}.

The main grid model uses a $4\times4\times32$ latent ($D=512$), a tile-transformer encoder, and a
factored spatial--temporal predictor with eight-frame history. Its representation training uses
neither pixels as prediction targets nor rewards. Each model is fitted on $500{,}000$ transitions
for approximately ten minutes on an A100; post-hoc decoders, heads, and policies are fitted separately.
The Craftax and gate checkpoints use the same grid shape and training family but different data.
Appendix~\ref{app:hparams} gives the complete settings and measured times.

\paragraph{Three tests.}
A frozen decoder tests whether information is \emph{recoverable from an encoding} under additional
supervision. Applying that same decoder to predicted latents tests whether it remains accessible in
\emph{model-generated futures}; because such latents can leave the encoder's distribution, a bad
rendering alone does not prove that the information is absent. Finally, a controller evaluated in
the real environment tests the \emph{complete learned pipeline}, including its observation interface,
reward and termination estimates, optimizer, and training distribution. These tests become more
informative when their shared components and changed components are explicit.

\paragraph{Reference model.}
For the gate, a reconstruction-trained tokenizer and causal flow-matching predictor provide a
feasibility reference \citep{lipman2023flowmatching,chen2024diffusionforcing,hafner2025dreamer4}.
It uses the same latent shape, data, and nominal model-training budget, with a 32-frame context and
jointly fitted reward and termination heads. Architecture, supervision, and history length differ
from the JEPA. One dynamics checkpoint is evaluated, so its policy runs show variation conditional
on that fit rather than variation across fitted flow models.

\section{Craftax: local goals and reward control ask different questions}
\label{sec:craftax}

Craftax-Classic combines movement, resources, inventory, and survival
\citep{matthews2024craftax}. It is a useful starting point because reaching a nearby visual target
and selecting a rewarding sequence of interactions can require different information. We evaluate
one fitted 512-dimensional grid JEPA in two protocols, both using real simulator outcomes.

\paragraph{Short-horizon goal reaching.}
After seven random warm-up actions, six additional random actions generate a reachable goal image.
The planner minimizes distance between its predicted endpoint latent and the encoded goal. It then
executes the selected six actions \emph{open loop}, without replanning. Success means matching the
goal's final player position; it does not require matching inventory or the entire scene. The planner
reaches 121 of 128 goals, compared with 24 for random actions (Table~\ref{tab:craftax}). Goal generation
and execution share simulator randomness. This is evidence for nearby goal reaching under the stated
protocol, not long-horizon task completion.

\paragraph{Reward-driven planning.}
We separately fit reward and value heads on frozen encodings and evaluate 100-step behavior from
128 starts. The six-step CEM controller replans after every real observation. Its mean return is
$0.47$, below $1.45$ for uniform random actions. A myopic reward controller, which does not use
predicted transitions, also underperforms random actions. The reward head's correlation falls from
$0.92$ on its fitting data to $0.35$ on held-out time rows, and value correlation falls from $0.99$
to $0.27$. Poor performance therefore already appears in reward-based action selection without predicted
transitions; the multi-step result concerns the combined head, dynamics, and planning pipeline.

\begin{table}[t]
\centering
\caption{\textbf{Different tasks, one Craftax checkpoint.} Goal reaching uses 128 nearby goals and
six open-loop actions. Reward control uses 128 starts and 100 steps; $\pm$ is the standard deviation
across starts. These are exploratory single-model results, not the standard Craftax achievement
benchmark.}
\label{tab:craftax}
\inputtable{craftax.tex}
\end{table}

This split motivates the controlled studies that follow. Goal-reaching success shows that a model
can guide some actions; it does not certify reward prediction or policy learning in that model.
Conversely, a weak reward pipeline does not establish that its encoder discarded the relevant
state. Appendix~\ref{app:craftax} records the planner and head-fitting details.

\section{The gate: information in the encoder and in imagination}
\label{sec:controllers}

\paragraph{A decision with a visible cause.}
A red cube moves in a $63\times63$ image. A vertical wall approaches from the right; its gap is at
the top or bottom with equal probability, but remains hidden until the wall passes a fixed reveal
point. From the centre, the gap becomes visible only two steps before contact, while either gap
requires five vertical moves. Retreating creates time to observe the gap and respond. Contact with
a wall or border ends the episode with reward $-1$; all other rewards are zero. Episodes last at
most 128 steps. The model-fitting data contain $500{,}000$ uniform-random transitions, with a separate
$25{,}000$-transition evaluation archive.

\paragraph{Comparisons and replication.}
We fit three independent grid-JEPA encoders. For each, recurrent PPO \citep{schulman2017ppo} is trained in the real gate
on frozen latents, and inside the frozen world model on predicted latents. The latent policies use
the same MLP torso design. In imagination, rewards and terminations come from fitted heads; in the
real gate they come from the simulator. Collection length, start states, update organization, and
step budgets also differ (Appendix~\ref{app:hparams}), making this a comparison of complete training
pipelines. Its real-versus-imagined contrast is not a dynamics-only intervention.

Each original predictor has three policy seeds per encoder, giving nine policies. A fitted
conditional residual predictor is evaluated with one sample and with a mean of eight samples,
adding nine policies per variant. Real-latent controls have nine policies; real-pixel controls have
three per torso. Every reported policy is its final checkpoint, evaluated on 256 real episodes.
The primary metric is walls passed per episode. We call survival above $50\%$ \emph{discovery} of
the survival strategy. Model fits and policy seeds are separate replication levels; the real and
imagined evaluations use different random-key schedules and are not paired trials.

\subsection{Useful wall information remains in the frozen representation}
\label{sec:results-wall}

Decoded real encodings visibly retain the approaching wall (Figure~\ref{fig:rollouts}). To test
whether that information can support behavior, we train a convolutional policy in the real gate
on the frozen decoder's output. Two of nine runs discover the survival strategy, both using encoder
seed 1; the corresponding real-pixel CNN discovers it in two of three runs
(Figure~\ref{fig:representation}). This establishes useful information in at least one fitted encoder
under the decoder-and-CNN read-out. The decoder adds reconstruction supervision outside the
world-model budget, so this is a test of recoverable information rather than unsupervised control
sufficiency.

The policy interface matters. With an MLP torso, real pixels yield $0.91$ walls per episode and
frozen latents yield $0.94$, with no discoveries in either group. An MLP on decoded latents similarly
yields $0.92$ walls. Thus failure of this MLP training recipe is not evidence that the encoder has
lost the obstacle. These comparisons do not set a capacity limit for all MLP policies; they locate
a limitation of the tested read-out and optimization procedure.

\begin{figure}[t]
\centering
\paperfig{rollouts.pdf}{1.0}
\caption{\textbf{A wall visible in real encodings disappears in predicted rollouts.} A pre-reveal
context is followed by futures under recorded actions: the simulator, decoded real JEPA encodings,
decoded original JEPA predictions, and one flow-model sample. The JEPA rows use encoder seed 0
and the same fitted decoder. The example is the first stored diagnostic context at wall position
$x=45$ among 64 contexts whose recorded futures survive the horizon. The flow sample is a single-checkpoint feasibility reference.
A decoding failure on predicted latents can reflect both prediction error and off-distribution
read-out; it is not alone proof of information loss.}
\label{fig:rollouts}
\end{figure}

\subsection{The tested imagination pipeline does not discover the strategy}
\label{sec:results-ppo}
\label{sec:results}

All three JEPA predictor variants produce poor downstream control
(Table~\ref{tab:gate-main}; per-seed outcomes in Appendix~\ref{app:round2}). The original predictor yields $0.27$ walls
per episode; one residual sample and the mean of eight each yield $0.26$. No policy among these
27 runs discovers the strategy; the highest final survival is $3.1\%$. Real-transition PPO on the
same frozen encoders yields $0.94$ walls. The difference measures the combined effect of learned
transitions, fitted heads, and the changed PPO protocol. Even the stronger real-latent MLP result
falls far short of reliable survival.

The flow reference supports feasibility at this task and nominal model budget: one of its nine
policies discovers the strategy. The remaining eight pass $0.72$--$1.18$ walls per episode. Since
all nine use one dynamics fit with a different architecture and supervision, this reference does
not establish a model-family ranking. Training the gate JEPA for twice and four times the steps
also leaves imagined-policy survival at $0$--$3\%$ in the tested extensions (encoder seed 0, policy
seed 0, each predictor variant).

\begin{table}[t]
\centering
\caption{\textbf{Final gate outcomes, with predictor variants kept separate.} Mean $\pm$ sample
standard deviation across policy runs; discovery is final survival above $50\%$. JEPA imagination variants and real-latent
conditions each contain $3$ models $\times$ $3$ policies; pixel controls have three policies. Real training changes rewards, terminations, and PPO protocol
alongside dynamics. Flow is a single-model reference.}
\label{tab:gate-main}
\inputtable{gate_main.tex}
\end{table}



\subsection{Prediction and fitted heads provide distinct failure signals}
\label{sec:results-mean}

The same decoder that renders real encoded walls removes them from original predicted rollouts.
Across all three encoder seeds, decoded shared-wall intensity is at most $0.02$ by the second
predicted step (Appendix~\ref{app:round2}). The shared wall pixels occur in both possible gap futures,
so a pixelwise mixture of those futures would retain them. This observation rules out that simple
pixel-space picture; it does not determine what averaging or distribution shift does in the learned
latent space.

A within-encoder counterfactual diagnostic makes the scale of the feature explicit. On 64 pre-reveal
contexts for seed 0, the first-step squared distance between wall-present and wall-absent encodings
is $0.035$, while total teacher-forced prediction error is $0.145$. The same ordering holds for all
three seeds (Appendix~\ref{app:round2}). This comparison concerns sensitivity to wall removal versus
error in all features. It neither identifies a causal mechanism nor permits comparing information
content through raw distances across independently fitted latent spaces.

The fitted heads expose another limitation. For encoder 0, termination AUC is $0.96$ for border
deaths but $0.65$ for wall deaths. Predicted rewards for the two death causes are $-0.89$ and $-0.05$,
respectively, against the true $-1$. An imagined policy therefore trains on a weak wall-collision
signal. The residual-sampling variants retain these heads and the policy protocol. Their poor control
outcomes therefore test this particular intervention, not distributional prediction in general.

\begin{figure}[t]
\centering
\paperfig{representation_by_seed.pdf}{1.0}
\caption{\textbf{Recoverable information can support successful control.} Final real-gate survival
for MLP policies on frozen latents or pixels (left), and CNN policies on decoded latents or pixels
(right). Points are policy runs, grouped by encoder seed; the dashed line is the $50\%$ discovery criterion. Both decoded-latent discoveries
use encoder 1. Reconstruction supervision and CNN training are additional to representation fitting.}
\label{fig:representation}
\end{figure}

\FloatBarrier
\section{Remove the wall: nuisance and dynamics interact}
\label{sec:cubes}

Cube worlds reduce the scene to one moving object. Random backgrounds either change independently
each frame or remain fixed for an episode; dynamics are deterministic or include hidden wind.
A fixed background is predictable but irrelevant to position. Per-frame noise is unpredictable and
also irrelevant. These cases distinguish visual complexity from predictive usefulness.

The strongest replicated comparison holds the 512-dimensional grid architecture, dataset size
($500{,}000$ transitions), and training length ($8{,}500$ updates) fixed across four gray-noise
conditions. A supervised nonlinear position read-out is fitted on training encodings and evaluated
on independent evaluation frames. It achieves subpixel error whenever either the background is
persistent or dynamics are deterministic. Combining per-frame noise with wind raises error to
$21.33\pm0.06$ pixels across three model initializations (Table~\ref{tab:cube-main}). This is a
failure to expose position through the tested read-out. Its prediction loss remains high, about
$0.73$ per coordinate, so it is not evidence of a successfully optimized low-loss nuisance solution.

\begin{figure}[t]
\centering
\paperfig{readout_by_seed.pdf}{1.0}
\caption{\textbf{The matched gray-noise cube comparison.} Position read-out from real encodings
for three independently initialized grid JEPAs in each condition. Architecture,
latent dimension, update count, and training dataset size are fixed. A high probe error is a failure
of the tested read-out, not proof of absent information.}
\label{fig:cube-readout}
\end{figure}

\begin{table}[t]
\centering
\caption{\textbf{A matched nuisance--dynamics interaction.} Grid-JEPA position L1 error in pixels,
mean $\pm$ sample standard deviation across three model fits. The same supervised read-out is
applied to real encodings and eight-step predictions; the latter can also reflect read-out
extrapolation. Predicted-window evaluation excludes resets and uses overlapping windows.}
\label{tab:cube-main}
\inputtable{cube_main.tex}
\end{table}

Other fitted models fail at different interfaces. On a black background without wind, the
class-token JEPA permits $0.087$-pixel position read-out from real encodings, but that read-out has
$45.27$-pixel error on eight-step predictions. A flow model also fails on persistent color texture,
with $18.56$-pixel real-latent error. These are single-fit observations with different dataset sizes,
reported separately in Appendix~\ref{app:cubes}; they preclude a uniform ranking of model families.
Together, the cube results motivate an objective-level question rather than a universal empirical
mechanism: what does low prediction loss plus non-collapse actually guarantee?

\section{What the objective does not guarantee}
\label{sec:theory}

The experiments concern fitted finite networks. The following population statements ask a narrower
question: which guarantees follow from prediction error and a Gaussian latent marginal alone?
They do not characterize optimization or identify the cause of the empirical failures.

\subsection{Non-collapse does not select the controlled state}
\label{sec:theory-rep}

Write observations as $x_t=R(p_t,b)$, where $p_t$ is controlled state and $b$ is an episode-constant
nuisance recoverable from the image. Assume that $b$ is independent of states and actions in the
training data. At deployment, assume it is independent of the initial state and transition randomness,
and does not directly affect dynamics or reward. Let $\mathrm{SIGReg}_\infty$ be the nonnegative
population characteristic-function discrepancy, which vanishes for $\Nrm(0,I_D)$
(Appendix~\ref{app:proofs}).

\begin{proposition}[Gaussian nuisance optimum; extending \citet{sobal2022slow}]
\label{prop:representation}
If the encoder class contains a map $f(R(p,b))=\varphi(b)$ with
$\varphi(b)\sim\Nrm(0,I_D)$ and the predictor class contains the last-latent copy map, their population
prediction loss and $\mathrm{SIGReg}_\infty$ are both zero. This is a global optimum with covariance
$I_D$ and $I(z_t;p_{1:T})=0$ under the training distribution. A deployed controller observing only
these latents, time, and its own actions has the value of a randomized state-blind policy.
\end{proposition}

The argument is simple: nuisance changes across episodes, satisfying the marginal constraint, but
is constant within each episode, making prediction exact. Its deployment value follows by giving a
state-blind policy the same private nuisance draw and controller randomness. The Gaussianization
assumption requires a rich nuisance and a realizable measurable map; a finite texture does not
literally instantiate it. The finite-batch regularizer also differs from the population functional.
A complementary innovation bound in Appendix~\ref{app:proofs} shows that an encoder preserving all
state distances at a fixed scale must incur prediction error for unpredictable state innovations.

\subsection{Decodability does not set a useful prediction scale}
\label{sec:amplitude}

Even preserving all state information leaves its geometry unconstrained. For any $\varepsilon>0$,
we can place the two values of a relevant bit less than $2\varepsilon$ apart while keeping the
\emph{entire} latent distribution exactly standard Gaussian.

\begin{proposition}[Exact information at arbitrarily small scale]
\label{prop:amplitude}
For every $\varepsilon>0$ and $D\ge1$, there is a measurable observation process with episode-constant
nuisance $b$ and a fresh independent fair state bit $p_t$, and an encoder $f_\varepsilon$, such that
$z_t\sim\Nrm(0,I_D)$ and both $p_t$ and $b$ are exactly recoverable almost surely. For fixed $b$,
$0<\|f_\varepsilon(R(1,b))-f_\varepsilon(R(0,b))\|<2\varepsilon$.
The optimal next-latent predictor is independent of the next bit and has coordinate-mean squared
error in $(0,\varepsilon^2/D]$. With reward $\mathbf1\{a_t=p_t\}$, latent feedback attains value $1$
per step, while every state-blind policy attains $1/2$.
\end{proposition}

The construction partitions one Gaussian coordinate into narrow intervals and uses the bit to
select the lower or upper half of each interval's conditional probability mass. The nuisance fixes
the interval and the within-half quantile. This preserves the Gaussian marginal and exact decoding
while making the two states arbitrarily close. The full proof is in Appendix~\ref{app:proofs}.
It uses infinite-precision measurable coding; the next bit is intrinsically unpredictable. It therefore
separates information from scale, without proving avoidable prediction error or explaining deletion
of the predictable gate wall.

\paragraph{Why scale matters for a downstream read-out.}
For a fixed $L$-Lipschitz value read-out $V$ and any coupled real and predicted next latents,
\begin{equation}
\label{eq:readout}
\left|\E V(Z')-\E V(\widehat Z')\right|
\le L\sqrt{\E\|Z'-\widehat Z'\|^2}.
\end{equation}
This standard bound follows from Lipschitz continuity and Cauchy--Schwarz. Prediction error controls
value error only relative to the read-out's sensitivity. Exact decodability alone supplies no bound
on that sensitivity. We do not estimate $L$ or claim a control guarantee for the fitted networks;
the inequality makes explicit the missing bridge between a geometric loss and a decision.

\section{Related work}
\label{sec:related}

\paragraph{Which JEPA setting?}
Image and video JEPAs learn by predicting representations
\citep{assran2023ijepa,bardes2024vjepa}. DINO-WM fits dynamics on frozen visual features, and V-JEPA~2
freezes its encoder for action-conditioned prediction with an L1 objective
\citep{zhou2024dinowm,assran2025vjepa2}. PLDM instead jointly learns its encoder and dynamics with
multiple regularizers \citep{sobal2025pldm}. Our direct setting is LeWorldModel's jointly trained
squared-loss predictor with SIGReg \citep{maes2026lewm,balestriero2025lejepa}; the population statements
do not automatically transfer to these other objectives.

\paragraph{Predictability, geometry, and planning.}
Predictable nuisance is an established limitation of joint embeddings \citep{sobal2022slow}.
Recent work studies how contrastive sampling changes slow-noise bias
\citep{gulati2026contrast}, and how temporal-distance objectives improve planning-relevant geometry
\citep{bai2026temporal}. LeWorldModel already reports accurate probes alongside weaker planning,
while \citet{terver2026planning} examine architecture and planner choices and discuss multimodal
conditional means. Self-predictive abstraction analyses similarly distinguish predictive structure
from task sufficiency \citep{tang2023selfpredictive,ni2024bridging}. Our contribution is a controlled
set of representation, rollout, and policy diagnostics with explicit replication, together with a
narrow separation between Gaussianity, exact decodability, and loss scale.

\paragraph{Models for decisions.}
Prediction--control mismatch predates JEPAs \citep{lambert2020objective}. Value-aware model learning,
value equivalence, and proper value equivalence ask which errors affect downstream Bellman
computations and policy values \citep{farahmand2017vaml,grimm2020value,grimm2021proper}.
Expectation models can suffice for suitable value-function classes \citep{wan2019expectation};
conditional-mean prediction is therefore not inherently incompatible with control. Stochastic
imagination models provide another route \citep{hafner2020dreamer,alonso2024diamond,hafner2025dreamer4},
but they too can fail at encoding, action conditioning, or rollout stages
\citep{hansen2026hallucination}. The same interface-level scrutiny should apply to both families.

\section{Discussion}
\label{sec:limitations}

The experiments separate capabilities that are often presented together. The Craftax checkpoint
supports nearby visual-goal planning, while its reward-planning pipeline has weak held-out reward and value
predictions. The gate encoders retain useful obstacle information under an added decoder and CNN,
while their tested imagination pipelines fail to discover the strategy. Nuisance cubes and the
population constructions show why a marginal non-collapse condition alone cannot settle either
question.

The practical consequence is to evaluate a world model at each interface: information recoverable
from real encodings, prediction of controlled counterfactuals, head generalization, and behavior in
the real environment. A controller's failure should lead to this decomposition before it is used to
judge the representation; a successful probe should be followed by prediction and control tests
before it is used to certify imagination.

Our empirical scope is one jointly trained squared-loss objective, small visual domains, and short
training budgets. Craftax and the flow reference each use one fitted model. Gate predictors, heads,
and PPO protocols have not been isolated by a fully matched intervention. The theory gives
population existence results, with measurable constructions rather than a finite-network or
optimization analysis. In particular, the fresh-bit example does not explain loss of a predictable
wall. What the evidence does establish is that non-collapse, decodability, and useful imagined
control are separate properties, requiring separate support.


\subsubsection*{Reproducibility statement}
Appendices~\ref{app:hparams}--\ref{app:craftax} specify model fitting, post-hoc supervision,
evaluation protocols, and replication levels. The tables retain final-checkpoint outcomes rather
than selecting intermediate peaks. The accompanying evidence manifest identifies the source
artifacts used for each result; independent policy seeds are distinguished from independent model
fits. Claims based on a single checkpoint are labeled throughout.

\subsubsection*{Statement on the use of AI}
Language models assisted with prose, code, plotting, proof checking, literature review, and
adversarial manuscript review. The authors are responsible for all claims, experiments, and
verification of the resulting text and artifacts.

\clearpage
\bibliography{references}
\bibliographystyle{iclr2027_conference}
\clearpage
\appendix
\section{Proofs}
\label{app:proofs}

\subsection{Empirical objective and population idealization}
For a training batch of $B$ windows and $H$ predicted time positions, the implemented loss is
\begin{equation}
\Lcal_{\mathrm{pred}}=\frac{1}{BHD}\sum_{b,h}\|\widehat z_{b,h+1}-z_{b,h+1}\|^2.
\end{equation}
The notation in Equation~\ref{eq:lewm} averages over time and uses this coordinate-mean convention.
One-step targets are teacher-forced at every temporal position, with a shorter history near the
start of a window. For encoded windows of $T$ frames, $P$ random unit projection directions $u_p$,
$K$ knots $\tau_k\in[0,3]$, and $w_k=e^{-\tau_k^2/2}$, the regularizer is
\begin{equation}
\label{eq:sigreg}
R=\frac{B}{TP}\sum_{t,p,k}q_k w_k
\left[\left(\frac1B\sum_b\cos(\tau_k u_p^\top z_{b,t})-w_k\right)^2
+\left(\frac1B\sum_b\sin(\tau_k u_p^\top z_{b,t})\right)^2\right],
\end{equation}
where $\Delta=3/(K-1)$, with $q_k=\Delta$ at the endpoints and $2\Delta$ in the interior.
The characteristic function is averaged across batch examples separately at each temporal position;
the factor $B$ is part of the implementation with $\lambda=0.09$.

The population statements use the nonnegative discrepancy
\begin{equation}
\label{eq:popsigreg}
\mathrm{SIGReg}_\infty(\mathcal L(z))=
\int_{\mathbb S^{D-1}}\!\int_0^3 w(\tau)
\left[\big(\E\cos(\tau\langle z,u\rangle)-w(\tau)\big)^2
+\big(\E\sin(\tau\langle z,u\rangle)\big)^2\right]d\tau\,d\sigma(u),
\end{equation}
where $w(\tau)=e^{-\tau^2/2}$ and $\sigma$ is uniform on the sphere.
Only the fact that this functional vanishes for a standard Gaussian is needed. It is not a theorem
about exact minima of the finite-batch objective.

\subsection{Nuisance optimum and the cost of innovations}
\paragraph{Proposition~\ref{prop:representation}.}
\begin{proof}
Let $\beta$ recover the nuisance from the observation, set $f^\star=\varphi\circ\beta$, and let
$g^\star$ copy the last latent. Then $z_t=\varphi(b)=z_{t+1}$, so every prediction residual vanishes.
The Gaussian marginal gives covariance $I_D$ and zero population regularizer. Both objective terms
are nonnegative, so their sum is globally minimized at zero. Training-distribution independence
gives $I(\varphi(b);p_{1:T})=0$.

For deployment, draw nuisance $B$ and the controller's random seed $U$ at reset, independently of
the initial state and transition randomness. Run the original controller in the environment with
nuisance $B$. Let a state-blind policy with private seed $(B,U)$ feed constant latent history
$\varphi(B)$ to a copy of that controller and reproduce its actions with $U$. Coupled to the same
initial state and transition randomness, the policies take the same first action. If their actions
agree through step $t$, their states, supplied latent histories, and seeds agree, so the next actions
agree as well. Induction gives identical trajectories, returns, and terminations. This does not
assume independence between $B$ and states subsequently influenced by its induced actions.
\end{proof}

\begin{corollary}[Innovation cost for uniformly state-resolving encoders]
\label{cor:strict}
Let $p_{t+1}=p_t+\mu(a_t)+w_t$, $\mathcal H_t=\sigma(p_{\le t},b,a_{\le t})$, and
$\Sigma_t=\Cov(w_t\mid\mathcal H_t)$. Assume finite second moments of states and encoded observations,
$x_t=R(p_t,b)$, fixed nuisance $b$, and
$\|f(R(p,b))-f(R(p',b))\|\ge\kappa\|p-p'\|$ for all admissible $p,p',b$, with $\kappa>0$.
Every predictor of latent and action history satisfies
\begin{equation}
\label{eq:innovation}
\frac1D\E\|g-f(x_{t+1})\|^2\ge\frac{\kappa^2}{D}\E\tr\Sigma_t.
\end{equation}
When $\E\tr\Sigma_t>0$, this exceeds the nuisance encoder's zero prediction loss under the
assumptions of Proposition~\ref{prop:representation}.
\end{corollary}
The assumption is stronger than injectivity or task sufficiency: every state distance must be
preserved at a uniform scale. Innovation is measured after conditioning on state history and nuisance;
persistent wind is not automatically fresh innovation.

\paragraph{Corollary~\ref{cor:strict}.}
\begin{proof}
Put $Z=f(x_{t+1})$ and $\mathcal C_t=\sigma(z_{\le t},a_{\le t})$.
Because observations and latents are measurable functions of state and nuisance, $\mathcal C_t\subseteq\mathcal H_t$.
Any predictor using a shorter window is also $\mathcal C_t$-measurable.
Conditional least squares and the tower property give
\begin{align*}
\E\|g-Z\|^2
&\ge\E\|Z-\E[Z\mid\mathcal C_t]\|^2\\
&=\E\|Z-\E[Z\mid\mathcal H_t]\|^2
 +\E\|\E[Z\mid\mathcal H_t]-\E[Z\mid\mathcal C_t]\|^2\\
&\ge\E\tr\Cov(Z\mid\mathcal H_t).
\end{align*}
The cross term vanishes because its first factor has conditional mean zero given $\mathcal H_t$.
Take conditionally independent copies $P,P'$ of $p_{t+1}$ given $\mathcal H_t$.
Such copies exist for Euclidean state spaces, using their regular conditional distributions.
The nuisance $b$ and current state $p_t$ are fixed under this conditioning.
The conditional independent-copy variance identity yields
\begin{align*}
\tr\Cov(Z\mid\mathcal H_t)
&=\tfrac12\E[\|f(R(P,b))-f(R(P',b))\|^2\mid\mathcal H_t]\\
&\ge\tfrac{\kappa^2}{2}\E[\|P-P'\|^2\mid\mathcal H_t]\\
&=\kappa^2\tr\Cov(p_{t+1}\mid\mathcal H_t)
=\kappa^2\tr\Sigma_t.
\end{align*}
The final equality follows because $p_t+\mu(a_t)$ is $\mathcal H_t$-measurable.
Taking expectations and dividing by $D$ proves Equation~\ref{eq:innovation}.
Averaging these inequalities gives the corresponding bound for a multi-time prediction objective.
Proposition~\ref{prop:representation} supplies the zero-loss comparator when its assumptions also hold.
\end{proof}

\subsection{Conditional prediction and geometric scale}
\begin{lemma}[Conditional mean]
\label{lem:mean}
For a fixed encoder and data law with $\E\|z'\|^2<\infty$, the unique squared-error minimizer
over measurable square-integrable predictors from history $c$ and action $a$ is
$m(c,a)=\E[z'\mid c,a]$, up to null sets.
\end{lemma}
\begin{proof}
Let $m=\E[z'\mid c,a]$.
For every square-integrable measurable $g(c,a)$, expand $g-z'=(g-m)+(m-z')$.
The cross term has expectation zero since $\E[m-z'\mid c,a]=0$.
Hence
\[
\E\|g-z'\|^2=\E\|g-m\|^2+\E\tr\Cov(z'\mid c,a).
\]
The second term does not depend on $g$.
The first is nonnegative and vanishes exactly when $g=m$ almost surely.
\end{proof}

\paragraph{The finite-sample regularizer.}
The finite-sample regularizer differs from Equation~\ref{eq:popsigreg}.
For independent identically distributed samples with characteristic function $\phi$ and empirical estimate $\widehat\phi_n$,
\begin{equation}
\E|\widehat\phi_n(t)-w(t)|^2
=|\phi(t)-w(t)|^2+\frac{1-|\phi(t)|^2}{n}.
\end{equation}
Indeed, $\E\widehat\phi_n=\phi$ and $\E|\widehat\phi_n-\phi|^2=(1-|\phi|^2)/n$, since each complex exponential has modulus one; expanding around $\phi$ gives the identity.
At a Gaussian law, the residual equals $(1-e^{-t^2})/n$.
This is estimator bias, not an optimization floor.
Changing the law changes both terms; dependent temporal samples introduce additional covariance terms. The implementation multiplies this statistic by batch size $B$, so its Gaussian finite-batch bias does not shrink as $1/B$ after that scaling.
Finite projections and knots also do not characterize Gaussianity.


\paragraph{Proposition~\ref{prop:amplitude}.}
\begin{proof}
Partition $\R$ into $I_j=[2\varepsilon j,2\varepsilon(j+1))$, $j\in\mathbb Z$.
Let $q_j=\Pr(G\in I_j)$ for $G\sim\Nrm(0,1)$.
Every $q_j$ is positive and their sum is one.
Let $F_j$ be the conditional Gaussian distribution function on $I_j$, with inverse quantile function $Q_j$ on $(0,1)$.
The conditional CDF is continuous and strictly increasing on the interval interior, so its inverse is measurable and takes values in that interior.

Draw $J$ with probabilities $q_j$, $U\sim\mathrm{Uniform}(0,1)$, and $V\sim\Nrm(0,I_{D-1})$, independently.
Set $b=(J,U,V)$; for $D=1$, omit $V$.
Take $R(p,b)=(p,b)$ and define
\begin{equation}
f_\varepsilon(R(p,b))=\big(Q_J((U+p)/2),V\big).
\end{equation}
The mixture $(U+p_t)/2$ is uniform on $(0,1)$, since $p_t$ is an independent fair bit.
Conditional on $J=j$, the first coordinate therefore has the Gaussian law restricted to $I_j$.
Mixing with weights $q_j$ recovers $\Nrm(0,1)$.
Independence of $V$ gives the joint law $\Nrm(0,I_D)$ and zero population regularizer.

Given the first coordinate $y$, its interval determines $J$ almost surely.
Write $v=F_J(y)$.
Then $p_t=\mathbf1\{v>1/2\}$ and $U=2v-p_t$ almost surely; the other coordinates recover $V$.
Thus the latent recovers both state and nuisance.
For fixed $b$, the two first coordinates are distinct and lie inside the same interval of width $2\varepsilon$.
All remaining coordinates agree, proving the distance bound.

Latent history reveals $b$, while the next bit remains independent and fair.
The next latent consequently has two equiprobable values, say $z^{(0)}(b)$ and $z^{(1)}(b)$.
Lemma~\ref{lem:mean} gives $m(b)=(z^{(0)}(b)+z^{(1)}(b))/2$.
Its conditional squared error is $\|z^{(1)}(b)-z^{(0)}(b)\|^2/4$, strictly positive and bounded by $\varepsilon^2$.
The midpoint is square-integrable by Jensen's inequality, since both state encodings have finite second moments.
Taking expectations and dividing by $D$ proves the coordinate-mean loss bound, including strict positivity since the conditional error is positive almost surely.
The midpoint depends only on $b$, so it is independent of $p_{t+1}$.
Finally, decoding the current bit permits $a_t=p_t$.
A state-blind action is independent of the current fair bit, giving success probability $1/2$.
\end{proof}


\section{The cube worlds}
\label{app:cubes}

The position-readout study crosses black, per-frame random, and per-episode random backgrounds
with deterministic or hidden-wind dynamics. Color-noise and black-background datasets contain
100,000 transitions; the gray-noise conditions contain 500,000. Only the four gray grid-JEPA
conditions have three independent model initializations. Table~\ref{tab:cube-main} and
Figure~\ref{fig:cube-readout} report that replicated comparison. Tables~\ref{tab:cube-readouts}
and~\ref{tab:cube-flow} retain seed-0 outcomes for the broader conditions, using the same canonical
benchmark files throughout. Differences across dataset sizes and model recipes are not isolated
nuisance effects.

\paragraph{Environment.}
The scene contains a red $6\times6$ object in a $63\times63$ RGB image with a red boundary. Four
cardinal actions move the object three pixels. Reset positions are uniformly sampled from integer
coordinates in $[16,42]^2$, and initial wind is zero. In wind conditions, each wind coordinate
independently increments by $-1$, $0$, or $1$ with probabilities $0.1$, $0.8$, and $0.1$, clipped
to $[-2,2]$; the updated wind is added to the action displacement. Background pixels are uniform
integers from 0 to 255; gray noise shares the same sample across RGB channels. Contact with the
boundary gives reward $-1$ and terminates; other rewards are zero, with a 128-step cap.

\paragraph{Read-out fitting.}
The nonlinear latent position probe is fitted on training data and evaluated on independent
\texttt{eval.npz} frames. Its spatial softmax-and-offset output is produced by two hidden layers of width 128, trained for
10,000 updates with batch 2,048 and Adam at initial learning rate $10^{-3}$ with cosine decay.
Grid inputs are standardized using training statistics. For class-token inputs, standardization and
PCA whitening are selected using the final 10\% of fitting rows, then the chosen variant is refitted
on the full fitting split. Whitening floors singular values at 1\% of the largest.
Applying this read-out
to predicted latents tests both prediction and read-out extrapolation. The pixel read-out used for
decoded flow images is a separate supervised CNN trained with synthetic background replacement
using ground-truth cube masks; it is a diagnostic, not a policy or a world-model training signal.
Rollout windows are sampled with replacement and overlap, and exclude death or reset during the
eight context and 16 forecast frames. The configured benchmark uses 2,048 windows for video metrics
and the first 512 for position diagnostics; exact sampled indices were not retained in the archived
result files. Surviving windows are a selected subset of wind trajectories, so these errors do not
establish death prediction or closed-loop control.



\begin{table}[t]
\centering
\caption{Seed-0 cube position L1 error (pixels) for the grid and class-token JEPAs. ``Real'' evaluates
held-out encoded frames; ``Pred. 8'' evaluates the same read-out on eight-step predicted latents.
Gray conditions use more training transitions than black and color conditions.}
\label{tab:cube-readouts}
\inputtable{cube_readouts.tex}
\end{table}

\begin{table}[t]
\centering
\caption{Seed-0 flow cube position L1 error (pixels). Latent and pixel columns use different
supervised read-outs. The persistent-color failure is visible before dynamics: real encoded position
has large error. These are model-recipe references rather than objective-only comparisons.}
\label{tab:cube-flow}
\inputtable{cube_flow.tex}
\end{table}

Earlier controller experiments use the three black/color background conditions. Their final
survival outcomes are retained below. The class-token JEPA's near-zero loss on per-episode texture
and poor reward-driven control are consistent with the known predictable-nuisance mechanism,
without exactly realizing Proposition~\ref{prop:representation}. The replicated gray-frame/wind
failure instead has high prediction loss and should not be interpreted as the same mechanism.
Still-world survival can be achieved by state-blind alternating actions, so it does not certify
state feedback.

\begin{table}[ht]
\centering
\caption{Survival (percent) in the supplementary cube worlds. ``LeWM'' denotes the class-token JEPA;
``grid'' denotes the grid JEPA; ``goal'' and ``heads'' use latent-goal and learned-return search,
respectively. PPO columns use the final saved policies; ``px'' and ``z'' identify pixel and latent
observations. Still-world survival does not certify state feedback.}
\label{tab:cube}
\inputtable{cube.tex}
\end{table}

\clearpage
\section{Hyperparameters and compute}
\label{app:hparams}

The gate and gray-noise cube conditions supply $500{,}000$ uniform-random transitions each. The earlier black and color-noise cube datasets contain $100{,}000$ transitions. Craftax uses its separate $500{,}000$-transition archive. Model training uses bfloat16 compute,
float32 parameters, and one A100. A $500$-step calibration estimates the training length fitting
$575$ seconds including compilation; the grid gate model trains for $8{,}500$ steps ($17{,}000$ and $34{,}000$ in the budget sentence of Section~\ref{sec:results}), the class-token cube models for $23{,}750$--$24{,}000$.
Model seeds $0$, $1$, $2$ of the grid JEPA differ only in the initialisation and data-order seed.

\paragraph{JEPA settings.}
Both JEPA architectures use batch $128$ and AdamW with learning rate and weight decay $10^{-3}$, $10\%$ warm-up
followed by cosine decay to zero, and gradient clipping at $1.0$. SIGReg uses weight $0.09$, $17$
knots on $[0,3]$, and $1024$ random projections. Neither architecture uses a stop-gradient, teacher, decoder, or
reward during world-model fitting. Linear--BatchNorm--GELU--linear projection heads flatten batch, time, and spatial tokens for normalization. Training-mode statistics therefore couple examples and temporal positions; at inference, running statistics define a fixed pointwise projection and the attention predictor is causal. The theoretical causal-map formulation is an inference/population idealization.

\begin{center}
{
\begin{tabular}{lll}
\toprule
& Class-token JEPA & Grid JEPA \\
\midrule
encoder            & ViT, $7\times7$ patches           & tile transformer $+$ conv \\
width / depth / heads & $192$ / $6$ / $3$              & $128$ / $6$ / $8$ \\
encoder MLP        & $768$                              & $4\times$ width \\
read-out           & class token                        & $4\times4$ grid, $32$ channels \\
latent dimension   & $192$                              & $512$ \\
projector          & BN projector, width $2048$      & same, hidden $512$ \\
predictor          & AdaLN-zero AR          & spatial + causal temporal \\
predictor depth / heads / dim & $6$ / $16$ / $64$  & $6$ / $8$ / $32$ \\
predictor MLP / width & $2048$ / latent                 & $2\times$ width / $256$ \\
history            & $3$ frames                         & $8$ frames \\
dropout            & $0.1$                              & $0$ \\
parameters         & $15.3$M                            & $2.6$M $+$ $8.3$M \\
\bottomrule
\end{tabular}}
\end{center}

\paragraph{Flow settings.}
The tokenizer uses $9{,}000$ steps, batch $256$, and learning rate $9\times10^{-4}$; dynamics
$16{,}000$ steps, batch $64$, the same learning rate, width $256$, six layers, eight spatial and
temporal attention heads, a $32$-frame context, independent per-frame noise levels, a shortcut loss of
weight $0.25$ every fourth step after the first quarter of training, history dropout $0.25$, two
history frames resampled every eight steps after half of training, and a temporal loss weight of
$0.25$. The symlog reward head has $255$ bins over $[-20,20]$; reward and termination losses have
weight $0.1$ and enter the trunk at full strength for half of training, then at factor $0.1$.
Inference uses four function evaluations per frame. We corrected terminal-transition masking in the
recipe's reward loss (Appendix~\ref{app:rewardbug}).

\paragraph{JEPA read-outs.}
Gate post-hoc heads use the frozen encoder and the first $800$ time rows ($80{,}000$ frames) of the same archive used for world-model training. The following $200$ time rows are held out from head fitting only; the independent world-model evaluation archive is \texttt{eval.npz}. The pixel decoder uses $8{,}000$ steps, batch $256$, learning
rate $10^{-3}$, width $256$, six residual blocks; it supports the figures and the decoded-latent control, in which its frozen output, clipped to $[0,255]$, replaces the pixels of the real-gate PPO; no true pixel enters that policy after the encoder. Reward,
termination, and behaviour-value heads are two-layer perceptrons of width $512$ on the latent and the
one-hot action, $6{,}000$ steps, batch $1024$, learning rate $10^{-3}$, Monte Carlo value targets at
discount $0.99$, value loss weight $0.1$. The residual predictor models $z_{t+1}-m$ with
$m=g(z_{t-7:t},a_{t-7:t})$ the original grid prediction, conditioned on $m$, $z_t$, and the action:
three residual blocks of width $1024$, $8{,}000$ steps, and batch $2048$. The final 20,000 valid
windows of the same world-model training archive are excluded from residual fitting and used for
its quality check. They are not independent world-model evaluation data. For encoder seed 0,
coordinate-mean MSE is $3.07\times10^{-4}$ for the original predictor, $2.05\times10^{-4}$ for one
residual sample, and $1.32\times10^{-4}$ for the mean of eight. This quality check uses eight function
evaluations per sample; the reported PPO runs instead use four. Averaging eight learned samples is
a random Monte Carlo average, not an exact conditional mean.

\paragraph{Policy optimisation.}
PPO uses discount $0.99$, GAE parameter $0.9$, clipping $0.2$, value coefficient $0.5$, entropy
coefficient $0.01$, AdamW at $2.5\times10^{-4}$ with weight decay $0.01$, four epochs per batch, and
gradient clipping at $1.0$. Pixels pass through a three-block Impala torso \citep{espeholt2018impala}, latents through a two-layer
perceptron; a GRU receives that embedding and the previous action; the critic uses $64$ categorical
bins with histogram-Gaussian targets \citep{farebrother2024stopregressing}. Real-environment PPO uses $128$ environments and $64$ steps per
iteration with four minibatch groups, truncated backpropagation over $16$ steps, and $1.5$M
environment steps. In-model PPO uses $150$ updates of $512$ trajectories of $16$ steps ($1.2$M
imagined steps) with minibatch size $128$; half of the starts are resets, half dataset states with
their real eight-frame context. All reported policy rows use final checkpoints.

\paragraph{Compute.} On a shared A100, the seed-0 grid-gate JEPA completes 8,500 updates in
573 seconds and has a logged total of 586.5 seconds. Probe and head fitting takes approximately
three minutes. The residual sampler's optimization stage takes 74.3 seconds; model loading, latent
extraction, fitting, and quality evaluation together take 275.2 seconds before the final file write.
Other logged stages are PPO inside the JEPA $2.0$--$4.3$ min per run; PPO inside the flow model
$5$--$6$ min per run at four function evaluations, $20$--$26$ min with the mean of eight; PPO in the real environment $3.7$--$7.7$ min on pixels with the convolutional torso and $51$--$70$ s with the MLP torso, $76$--$78$ s on the frozen latent, $4.7$--$4.8$ min on the decoded latent with the convolutional torso and $1.8$--$2.9$ min with the MLP torso; the latent diagnostic: $2$ min.

\section{Evaluation protocol}
\label{app:protocol}

Every gate-policy endpoint uses the final checkpoint and $256$ real episodes from reset, with the
$128$-step cap. Survival is the fraction reaching the cap. Walls passed counts walls leaving the
frame while the agent remains alive. Discovery is final survival strictly above $50\%$.
Real-training and imagined-training runs advance their random keys under different iteration schedules; their final evaluations are not paired episode trials. Policy seeds are separate
runs conditional on a fitted model; independently fitted encoders provide a different replication
level. The flow model has only one checkpoint. Reported means weight policy runs equally;
standard deviations describe variation across those runs, not uncertainty across model families.

\paragraph{Gate specification.}
The gate uses the same image and object sizes and cardinal action displacement. A white vertical
wall of width three pixels spawns at $x=60$ every 22 steps and moves left three pixels per step.
Its 16-pixel gap is sampled at spawn, top or bottom with equal probability, and becomes visible
when the wall reaches $x\le40$. The gap location is hidden state before that reveal. Border or wall
contact produces reward $-1$ and termination; the time cap is 128 steps.

\paragraph{Latent diagnostic.} For each encoder, the diagnostic uses $64$ real contexts with the
wall still solid at $x\in\{42,45\}$. Reconstructed simulator states generate recorded futures and
matched wall-free futures under the recorded actions. The recorded branch survives the diagnostic
horizon by construction. At each step, wall amplitude is the squared Euclidean distance between
encodings of the wall-present and wall-absent frames. The teacher-forced residual uses real encoded
history and is a total squared error, summed over latent coordinates.
The wall coordinate of a prediction is its projection onto the wall-present minus wall-absent
latent difference, with endpoints $0$ and $1$. Decoded wall intensity averages the minimum colour
channel on shared wall pixels, excluding either possible gap band. These measurements compare
features within one encoder; raw distances across models are not an information metric.

\section{Latent-policy torso controls}
\label{app:latentctl}

The MLP on decoded latents passes $0.92$ walls per episode with no discoveries in six runs
(Table~\ref{tab:latentctl}), using encoders $0$ and $1$. The decoded CNN control instead uses all
three encoders, with three policies per encoder. Both of its discoveries use encoder $1$.
The decoder is frozen and reconstruction-supervised; success therefore establishes useful encoded
information under this added read-out, not sufficiency of the unsupervised training objective.
Failure of the true-pixel MLP recipe establishes no upper bound on MLP capacity.

\begin{table}[ht]
\centering
\caption{Real-gate observation and torso controls: final walls per episode and survival, mean
$\pm$ standard deviation over policies; discovery means final survival above $50\%$.
All policies use $1{,}499{,}136$ real transitions.}
\label{tab:latentctl}
\inputtable{latentctl.tex}
\end{table}

\section{Control outcomes and latent diagnostics}
\label{app:round2}

Table~\ref{tab:paired} gives aggregate final outcomes; Table~\ref{tab:seeds} lists every
imagined JEPA policy. The pooled imagined JEPA mean weights all original, sampled-residual and
mean-residual runs equally. The flow rows distinguish six sampled policies from three policies
using a mean of eight, all sharing the same model fit. They do not estimate variation across
independent flow-model fits. Tables~\ref{tab:seeddiag} and \ref{tab:latent} retain the encoder
and geometry diagnostics supporting the gate observations.

\begin{figure}[t]
\centering
\paperfig{control_by_seed.pdf}{1.0}
\caption{\textbf{Control outcomes by model fit and predictor variant.} Each point is a final policy
evaluated on 256 real episodes. Original prediction, residual sampling, and residual means are
separate conditions; model and policy replication are distinct. Horizontal bars show within-model means,
and the dashed line marks random actions. The flow reference uses one checkpoint.}
\label{fig:control}
\end{figure}

\begin{table}[ht]
\centering
\caption{Final control outcomes in the real gate, mean $\pm$ standard deviation over policy runs.
Imagined training uses $1{,}228{,}800$ transitions per policy; real training uses $1{,}499{,}136$.
Model-fit counts and policy-run counts identify distinct replication levels.}
\label{tab:paired}
\inputtable{paired.tex}
\end{table}

\begin{table}[ht]
\centering
\caption{Every imagined JEPA policy at its final checkpoint: survival (percent) and walls per
episode by encoder seed, predictor variant and PPO seed. Each cell uses $256$ real episodes.}
\label{tab:seeds}
\inputtable{seeds.tex}
\end{table}

\begin{table}[ht]
\centering
\caption{Grid-JEPA diagnostics by encoder seed. Wall amplitude and total teacher-forced residual
are squared distances at future steps $1$ and $8$; decoded intensity uses steps $1$ and $2$.
The wall coordinate has endpoints $0$ (wall absent) and $1$ (true encoding).
Head AUC compares each death cause with alive transitions; PPO wall counts pool predictor variants.}
\label{tab:seeddiag}
\inputtable{seeddiag.tex}
\end{table}

\begin{table}[ht]
\centering
\caption{JEPA geometry and decoded wall intensity for model seed $0$, averaged over $64$
pre-reveal contexts. Distances use matched actions and the same encoder; total residuals include
all latent coordinates. Wall intensity is measured on shared wall pixels.}
\label{tab:latent}
\inputtable{latent.tex}
\end{table}

\begin{figure}[t]
\centering
\paperfig{geometry_by_seed.pdf}{1.0}
\caption{\textbf{Within-encoder scale diagnostics.} For each of the three gate encoders, wall-removal
distance and teacher-forced total one-step residual are summed squared latent distances. Curves
average 64 pre-reveal contexts whose recorded futures survive the diagnostic horizon. These are
feature sensitivity and total prediction error, not commensurate measures of information content.}
\label{fig:geometry}
\end{figure}

\clearpage
\section{Terminal rewards in the flow recipe}
\label{app:rewardbug}

The original recipe masks transitions entering reset frames out of reward training. Here the only
nonzero reward occurs on terminal transitions, so that mask removes the reward signal.
The corrected head predicts the next reward from the current frame summary and action with
terminal transitions retained in its loss. The reported flow-policy outcomes use this corrected
convention. The single-checkpoint reference therefore denotes the complete tokenizer, dynamics,
reward and termination package, rather than an isolated change in transition stochasticity.


\section{Craftax protocols and single-checkpoint results}
\label{app:craftax}

The grid checkpoint uses training seed 0, a tile encoder, $16\times32=512$ latent coordinates,
a six-layer factored predictor of width 256, history length eight, 17 actions, batch 128, and
8,500 updates. The recorded parameter count is 11.12 million. Training takes 604.5 seconds from
loading through the logged end-to-end timer (578.5 seconds in the update loop); thus it is an
approximately ten-minute run, not a strict sub-600-second result. The frozen decoder and heads add
175.6 seconds. Results are from one model fit and evaluation seed 0.

\paragraph{Goal generation and planning.}
Reset states receive seven random warm-up actions. Six further random actions generate a reachable
goal image. CEM optimizes a length-six action sequence in latent space, with population 512,
32 elites, and six iterations. The learned planner uses the goal image, without privileged position
coordinates or the generating actions in its cost. It executes the entire selected sequence without
observation feedback. Success is equality of final player positions. Goal generation, random actions,
and the selected plan use the same future simulator random keys; replaying the generating actions
therefore reaches all goals and serves only as a protocol check. This protocol samples nearby
reachable goals, not arbitrary distant tasks.

The planned endpoint position L1 error is $0.0781$, against $1.5313$ for random actions. Pixel MSE
on the 0--255 scale is $104.00$ versus $1477.92$. These quantities use the same 128 starts as the
reported $121/128$ versus $24/128$ reach counts. They do not quantify variation across fitted models.

\paragraph{Reward control and head generalization.}
Reward and value heads use the first 80,000 transitions (320 time rows across 250 streams) from the
world-model training archive. Held-out time rows are excluded from head fitting but not from encoder
training. Value fitting uses the first 80 time rows owing to truncated-return handling. Reward and
value fitting correlations are $0.922$ and $0.995$; held-out correlations are $0.353$ and $0.266$.
Held-out reward RMSE is $0.222$. Mean predicted held-out value is $2.356$, versus mean empirical
discounted return $1.348$.

Reward evaluation uses 128 starts, 100 steps, no autoreset, and rewards accumulated while alive.
The CEM controller executes one action and replans from each real observation. Random actions,
myopic fitted-reward argmax, one predicted step plus value, and CEM obtain returns
$1.4453$, $0.7813$, $0.9438$, and $0.4734$, respectively; their across-start standard deviations
are $1.0897$, $1.1109$, $1.2014$, and $1.1039$. Fractions alive at the end are $0.6328$, $0.7266$,
$0.6328$, and $0.6719$. Aggregate summaries do not preserve the per-start covariance needed for
paired confidence intervals. The archive's mean reward multiplied by 100 is not an online behavior
baseline and is excluded from the comparison.

Earlier Craftax experiments used a different $81\times32$ spatial latent and different horizons.
They are not merged with these 512-dimensional results. There is no matched online flow-controller
comparison for this checkpoint.

\end{document}
